Este anexo describe los pasos y herramientas necesarias para poder ejecutar la aplicación a partir de las imágenes publicadas en DockerHub. Dado que se está trabajando con imágenes que se comunican a través de un Docker Compose, será necesario tener Docker desktop instalado en ordenadores Windows y Mac, o Docker y Docker Compose en ordenadores Linux. 

Los pasos para ejecutar la aplicación son los siguientes:
\begin{enumerate}
    \item Descargar el motor de Docker y la herramienta Docker Compose. Para ello, en ordenadores Windows y Mac, es necesario descargar la aplicación Docker Desktop desde su \href{https://www.docker.com/products/docker-desktop/}{\uline{página principal}\footnote{\url{https://www.docker.com/products/docker-desktop/}}}. Para ordenadores Linux, se deben seguir las instrucciones incluidas en la documentación para instalar el motor de Docker y la herramienta Docker Compose. Instrucciones de instalación del \href{https://docs.docker.com/engine/install/ubuntu/}{\uline{motor de Docker}\footnote{\url{https://docs.docker.com/engine/install/ubuntu/}}}. Instrucciones de instalación de \href{https://docs.docker.com/compose/install/linux/}{\uline{Docker Compose}\footnote{\url{https://docs.docker.com/compose/install/linux/}}}
    \item Descargar el archivo docker-compose.yml del \href{https://github.com/codeurjc-students/2024-Atra}{\uline{repositorio de la aplicación}\footnote{\url{https://github.com/codeurjc-students/2024-Atra}}}.
    \item Asegurarse que el motor de Docker está en ejecución. En Windows y Mac, con Docker Desktop abierto, en la esquina inferior izquierda indica el estado del motor. Debe estar en ejecución. En Linux, podemos comprobar el estado ejecutando \texttt{sudo systemctl status docker}. Si no está activo, debemos iniciarlo con \texttt{sudo systemctl start docker}.
    \item Con el motor en ejecución, podremos lanzar la aplicación ejecutando el comando docker-compose up.
    \item La aplicación está en ejecución. Podemos acceder a ella a través de localhost:80. La base de datos está disponible en localhost:3307, podemos utilizar herramientas como MySQL Workbench para conectarnos a la base de datos y monitorizar su estado. El \textit{backend} está disponible en http://localhost:8080, y se puede interactuar con él mediante herramientas como postman.
\end{enumerate}

Una vez ejecutada la aplicación, podremos iniciar sesión con una de las cuentas creadas, o crear una nueva. En los datos de ejemplo que se cargan por defecto existen tres usuarios: "angel", "david", y "ruben". Las contraseñas de cada usuario son el nombre de usuario. Cada una de esas cuentas tiene una serie de actividades y rutas definidas. La siguiente tabla describe las rutas, indicando su visibilidad, a qué usuario pertenecen, y cuantas actividades de cada usuario hay en cada ruta:
\begin{table}[H]
\centering
\begin{tabular}{l l c c c}
\hline
\textbf{Ruta} & \textbf{Visibilidad} & \textbf{Angel} & \textbf{David} & \textbf{Ruben} \\
\hline
Laviana circular corta & Pública & 2 & 2 & 3* \\
Laviana circular larga & Mural Specific & 0 & 2 & 3* \\
Laviana lineal & Privada & 0 & 0 & 3* \\
Mostoles 10k & Mural Specific & 5* & 2 & 1 \\
Hospital - Soto & Pública & 3* & 4 & 0 \\
Soto Vuelta & Privada & 2* & 0 & 0 \\
Parque Agustin & Pública & 2 & 4* & 0 \\
Forma L & Mural Specific* & 0 & 2* & 0 \\
Pista Valle & Privada & 0 & 1* & 0 \\
\hline
\end{tabular}
\caption{Listado de rutas y actividades por usuario}
\end{table}
Algunas anotaciones adicionales sobre los datos de ejemplo:
\begin{itemize}
    \item Los asteriscos en el número de actividades en la ruta indican el propietario de la ruta.
    \item La ruta "Forma L" tiene su visibilidad marcada con un asterisco ya que, aunque es mural specific, no especifica ningún mural, por lo que es funcionalmente privada.
    \item Existe un mural creado por el usuario "angel". Todos los miembros forman parte del mural.
    \item Existe un usuario adicional, "extra". Este usuario no tiene actividades ni rutas. Simplemente es propietario de una serie de murales, para que aparezcan murales en la sección discover y los usuarios reales se puedan unir a ellos. Este usuario no forma parte del mural creado por "angel", y los otros usuarios no forman parte de ninguno de los murales creados por este usuario.
\end{itemize}

Adicionalmente, existen otros dos sets de datos de ejemplo, aunque para acceder a ellos será necesario modificar el docker-compose.yml. Concretamente, se deberá modificar la variable de entorno \texttt{INIT\_DB} de la sección \texttt{environment} del servicio \texttt{backend}. Por defecto está fijado a "demo2", que corresponde a los datos descritos. Adicionalmente existen:
\begin{itemize}
    \item \textbf{demo1}: inicializa la aplicación con un único usuario "qwe", que contiene una serie de actividades y rutas.
    \item \textbf{demo3}: inicializa la aplicación con tres usuarios, asd, qwe, y zcx. Cada uno tiene una serie de actividades, rutas y murales con distintas visibilidades. Sin embargo, las agrupaciones realizadas son totalmente arbitrarias; se agrupa en una misma ruta actividades que ocurren en rutas distintas. Es el dataset utilizado durante el desarrollo.
    \item \textbf{off}: desactiva la inicialización, de forma que la aplicación carga vacía.
\end{itemize}

